===== BEGIN FILE: Short_arxiv_submission.tex =====


\documentclass[
 floatfix,twocolumn,superscriptaddress, prl,aps










 amsmath,amssymb,
 aps,






]{revtex4-2}
\usepackage[utf8]{inputenc}
\usepackage{amsmath}
\usepackage{amsfonts}
\usepackage{amssymb,amsthm}




\usepackage[subrefformat=parens,labelformat=parens,caption=false]{subfig}

\usepackage{graphicx}
\usepackage{multirow}
\usepackage[]{hyperref}
\usepackage{bbold}
\usepackage{cancel}
\hypersetup{
  colorlinks = true,
  urlcolor = magenta,
  linkcolor = red,
  citecolor = blue
}
\usepackage[normalem]{ulem}
\usepackage{braket}
\usepackage[all]{hypcap}
\usepackage{url}
\usepackage{float}








\usepackage{xcolor}


\newcommand{\red}[1]{\textcolor{red}{#1}}
\newcommand{\blue}[1]{\textcolor{blue}{#1}}
\newcommand{\orange}[1]{\textcolor{orange}{#1}}
\newcommand{\brown}[1]{\textcolor{brown}{#1}}
\newtheorem{theorem}{Theorem}
\newtheorem{lemma}[theorem]{Lemma}
\newtheorem{corollary}{Corollary}[theorem]

\usepackage[]{hyperref}
\usepackage{bbold}
\usepackage{bm}
\usepackage{cancel}
\newcommand{\eq}[1]{\begin{align}#1\end{align}}
\newcommand{\ie}{i.\,e.,\ }







\usepackage{graphicx}
\graphicspath {{figures/}}

\begin{document}

\title{Universal two-excitation scattering in two-dimensional subwavelength atomic arrays
}

\author{Yidan Wang}
\affiliation{Department of Physics, Harvard University, Cambridge, Massachusetts 02138, USA}
\author{Oriol Rubies-Bigorda}
\affiliation{Physics Department, Massachusetts Institute of Technology, Cambridge, Massachusetts 02139, USA}
\author{Valentin Walther}
 \affiliation{Department of Physics and Astronomy, Purdue University, West Lafayette, Indiana 47907, USA}
 \affiliation{Department of Chemistry, Purdue University, West Lafayette, Indiana 47907, USA}

\author{Susanne F. Yelin}
\affiliation{Department of Physics, Harvard University, Cambridge, Massachusetts 02138, USA}

\begin{abstract}
Subwavelength atomic arrays are a leading platform for engineering strong light-matter interactions, presenting exciting opportunities for quantum science. However, a full understanding of their multi-excitation dynamics remains a significant challenge. In this work, we uncover a remarkable universal phenomenon that emerges in these arrays. Using scattering theory to analyze two-excitation interactions, we reveal a profound simplification near critical points of the collective atomic excitation band structure, determined solely from single-excitation properties. At these critical points, scattering becomes universal
 and the full two-excitation scattering matrix decomposes into a block-diagonal form.

Remarkably, all scattering processes involving the photon field are completely suppressed, resulting in the perfect isolation of a unitary, nonlinear interaction channel between collective dark spin waves. Our findings provide exact analytical insights into few-body nonlinearities and establish a universal framework for analyzing complex scattering phenomena in ordered atomic systems.


\end{abstract}
\maketitle




\begin{figure}
    \centering
\includegraphics[width=1\linewidth]{illu_2D_c.pdf}

\caption{
   \textbf{(a)} Schematic of a subwavelength atomic array in real space, where the lattice spacing $d$ is smaller than half the resonant atomic wavelength ($\lambda_0$).
   \textbf{(b)} Universal structure of the S-matrix at critical energies. The S-matrix, $S_{\alpha\beta}$, connects input ($\alpha$) and output ($\beta$) channels defined by their constituent excitations, shown in momentum space. These excitations are photons and dark spin waves. Each photon (orange arrow) has an in-plane momentum corresponding to a bright spin wave (dashed arrow) inside the light cone (LC), while dark spin waves (blue wavy arrows) have momenta outside the LC. Generically, all kinematically allowed channels are coupled (faint gray arrows). At the critical energies, however, the S-matrix becomes block-diagonal: the purely atomic channel ($\alpha,\beta=0$) becomes unitary, while all nonlinearities are suppressed in the photonic sectors ($\alpha,\beta > 0$), leaving only single-particle dynamics.
}
    \label{fig:scattering_illu}
\end{figure}

 The pursuit of universal principles that distill simple, predictable behaviors from complex underlying dynamics is a central goal of modern physics. This principle of universality reveals that disparate systems can exhibit identical limiting behaviors, irrespective of their microscopic intricacies. For instance, in condensed matter, different physical systems can exhibit identical critical exponents near a phase transition, determined only by fundamental symmetries rather than microscopic details \cite{Goldenfeld1992, Chaikin1995}. Similarly, in fluid dynamics, the Navier-Stokes equations describe  the long-wavelength behavior of materials as different as water and plasma using only a few universal parameters like viscosity and density \cite{Landau1987, Batchelor1967}.

Here, we demonstrate that this principle of universality provides a powerful framework for understanding the nonlinear quantum  optics of subwavelength atomic arrays. These systems, where atoms are spaced closer than their resonant wavelength, have emerged as promising platforms to engineer light-matter interactions and advance various quantum technologies~\cite{bettles2016enhanced, ExponentialAsenjoGarcia2017, shahmoon2017cooperative, Manzoni2018, bekenstein2020quantum, Patti2021, rui2020subradiant, wei2021generation, srakaew2023subwavelength, Shah2024, Rubies-Bigorda2025Deterministic}. In the single-excitation limit, their collective behavior is well understood: they act as tunable mirrors for photons and support long-lived, dark spin waves~\cite{Manzoni2018, ExponentialAsenjoGarcia2017,shahmoon2017cooperative, bettles2016enhanced}.
Together, photons and spin waves serve as the fundamental excitations of the system.

However, harnessing atomic arrays for quantum nonlinear optics requires moving beyond this single-excitation picture. Two-excitation scattering is a crucial step into this nonlinear regime, but it presents a formidable challenge. A complete description must track the many possible scattering processes between incident pairs of photons and dark spin waves (see Fig.\ \ref{fig:scattering_illu}), a problem complicated by the mixed dimensionality of the excitations (3D photons and 2D spin waves). Consequently, existing approaches often rely on numerical simulations, which are inherently limited to finite systems~\cite{moreno2021quantum, zhang2022photon}, or employ various approximations~\cite{cidrim2020photon, bettles2020quantum, williamson2020superatom, Ostermann2024, RubiesBigorda2023, Scarlatella2024, scarlatella2024subwavelength}. Other methods are restricted to specific scenarios ~\cite{Solomons2023, Pedersen2024, Zhang2019, Iversen2021}, leaving a holistic, analytical understanding elusive.

In this Letter, we overcome this complexity by discovering a universal structure governing two-excitation scattering. We demonstrate that near specific characteristic energies, the entire two-body S-matrix (which encompasses all photon-photon, photon-spin wave, and spin wave-spin wave processes) adopts a universal form. Crucially, the energies where this universality emerges can be predicted from the system’s single-excitation properties alone, circumventing the need to solve the full, complex two-body problem. We employ a comprehensive scattering formalism, developed in a parallel work~\cite{Wang2025MultiExcitation}, to calculate the relevant scattering observables and thereby analytically verify this universal behavior. This discovery offers a powerful predictive shortcut, providing holistic and analytical insights into the rich nonlinear scattering properties of 2D atomic arrays.



\textit{System Hamiltonian.---}The total Hamiltonian $H = H_{\text{quad}} + U$, derived in detail in our parallel work \cite{Wang2025MultiExcitation}, comprises a single-excitation part $H_{\text{quad}}$ and a two-excitation interaction $U$. $H_{\text{quad}}$ conserves lattice momentum $\bm{p}$ (the quasi-momentum along the array dimension) and decomposes as
$H_{\text{quad}} = \int_{\text{BZ}} d\bm{p} \, H_{\bm{p}}$~\footnote{Here and throughout, $d\bm{k}$ denotes the integration measure over the 2D momentum space, i.e., $d^2k = dk_x dk_y$.}, where BZ denotes the atomic Brillouin zone.
 Each component $H_{\bm{p}}$ governs the dynamics of collective atomic spin waves, created by the bosonic operators $b^\dagger_{\bm{p}}$, and their coupling to the corresponding photon channels.
 The atomic resonance frequency $\omega_{eg}$ defines the light cone radius $k_0 = \omega_{eg}/c$. In the subwavelength regime, where the lattice spacing $d < \lambda_0/\sqrt{2}$ (with $\lambda_0 = 2\pi c/\omega_{eg}$), this light cone is contained entirely within the Brillouin zone, thereby defining two physically distinct regions of momentum space. We work in a rotating frame at the atomic resonance frequency $\omega_{eg}$, effectively setting it as the zero of energy for excitations.
Outside the light cone ($\|\bm{p}\| > p_0$), excitations $b^\dagger_{\bm{p}}$ are decoupled from propagating photons, forming long-lived dark spin waves with Hamiltonian $H_{\bm{p}} = \Delta(\boldsymbol{p}) b^\dagger_{\boldsymbol{p}} b_{\boldsymbol{p}}$, where $\Delta(\bm{p})$ is the  collective Lamb shift obtained by tracing out off-resonant vacuum electromagnetic modes~\cite{ExponentialAsenjoGarcia2017}.
Inside the light cone ($\|\bm{p}\| < p_0$), excitations $b^\dagger_{\bm{p}}$ couple radiatively to effective 1D photon channels $C_{\bm{p}}(\chi)$ and are referred to as bright spin waves. Each channel $C_{\bm{p}}(\chi)$ represents photons propagating away from the array with 3D wavevectors $(\bm{p}, \chi)$ and $(\bm{p}, -\chi)$, sharing the same in-plane momentum $\bm{p}$ as the spin wave $b_{\bm{p}}$. Here, $\chi( >0)$ is the magnitude of the out-of-plane momentum. The Hamiltonian in the light cone reads
\begin{align}
H_{\bm{p}} ={}&   \Delta(\boldsymbol{p}) b^\dagger_{\boldsymbol{p}} b_{\boldsymbol{p}} + \int_{0}^{\infty} \!\!\! d\chi \, E_{\bm{p}}(\chi) C_{\boldsymbol{p}}^{\dagger}(\chi) C_{\boldsymbol{p}}(\chi) \nonumber \\
&+ \int_{0}^{\infty} \!\!\! d\chi \, \left[ g_{\boldsymbol{p}}(\chi) C_{\boldsymbol{p}}^{\dagger}(\chi) b_{\boldsymbol{p}} + \text{h.c.} \right] \quad (\text{for } \|\bm{p}\| < p_0),
\label{eq:Hp_bright_rev}
\end{align}
with photon energy $E_{\bm{p}}(\chi) = c\sqrt{\|\bm{p}\|^2 + \chi^2}-\omega_{eg}$. The dispersion $\Delta(\bm{p})$ for lattice constant $d=0.2\lambda_0$ is shown in Fig.~\ref{fig:results_all}(a).

The effective interaction $U$, resulting from the two-level nonlinearity of atoms ($u \to \infty$), conserves total momentum $\bm{P}$ while scattering relative momentum $\bm{q} \to \bm{q}'$:
\begin{equation}
U = \frac{u}{2} \! \int \! d\bm{q} \, d\bm{q}' \! \int \! d\bm{P} \, b^\dagger_{\boldsymbol{P}/2+\bm{q}'} b^\dagger_{\boldsymbol{P}/2-\bm{q}'} b_{\boldsymbol{P}/2-\bm{q}} b_{\boldsymbol{P}/2+\bm{q}}.
\label{eq:U_interaction_revised}
\end{equation}
This term mediates scattering between atomic spin waves.




\emph{Two-excitation scattering.---} We focus on the two-excitation manifold where scattering involves pairs of photons (\(C^\dagger_{\bm{p}}(\chi)\), \(\|\bm{p}\| < p_0\)) and/or dark spin waves (\(b^\dagger_{\bm{p}}\), \(\|\bm{p}\| > p_0\)), which serve as the stable asymptotic states of the  scattering process. Bright spin waves (\(b^\dagger_{\bm{p}}, \|\bm{p}\| < p_0\)) are transient states that decay into photons, making them unstable and unsuitable as asymptotic states. Incoming photons can be prepared using laser beams, while dark spin waves can be initialized through techniques such as rapid modulation of the lattice spacing, dynamic light shifts \cite{Rubies-Bigorda2022}, or multi-photon excitation schemes \cite{Rusconi2021}.




 \begin{figure}
     \centering     \includegraphics[width=1\linewidth]{result_figure4.pdf}
\caption{
Universal scattering signatures near critical points for \(\bm{P}=\bm{0}\).
We consider a square array with lattice spacing \(d = 0.2\lambda_0\) and \(\sigma^+\)atomic polarization with respect to the out-of-plane quantization axis.

\textbf{(a)} The two-dark-wave dispersion relation \(\Delta^{(2)}(\bm{0}, \bm{q})=2\Delta(\bm{q})\). The blue box outlines the two-excitation Brillouin zone, \(BZ^{(2)}\). Red and yellow markers correspond to the band maximum (\(E_{\text{max}}\)) and saddle points (\(E_{\text{sadd}}\)), respectively. The orange semicircle (\(D_2\)) marks the subspace of two bright spin waves.

\textbf{(b)} The non-trivial eigenvalue of the S-matrix for two incoming dark spin waves, $s(E)$, as a function of their total energy $E$. In agreement with the universal prediction of Eq.~\eqref{eqSEmax0_final}, the squared magnitude \(|s(E)|^2\) (black line) approaches 1 at the critical points, while the phase \(\text{arg}[s]\) (green dashed line) approaches 0 (at \(E_{\text{max}}\)) and \(\pm\pi\) (at \(E_{\text{sadd}}\)).
\textbf{(c)} Cross section for incident dark waves (\(\alpha=0\)) as a function of their incoming relative momentum \(\bm{q}\). \textbf{(c1)} The total cross section \(\sigma_{0,\text{tot}}\) is enhanced at the critical points and suppressed along the saddle lines. \textbf{(c2)} The branching ratio into photons, \(\sigma_{0,2}/\sigma_{0,\text{tot}}\), vanishes at the saddle points and along the saddle lines, but remains finite at the maximum, revealing the distinct universal character of the two critical energies.

\textbf{(d)} For two normally incident photons (\(\alpha=2\)), all cross sections vanish at the critical energies, confirming the decoupling of photons and dark spin waves.
}
\label{fig:results_all}
 \end{figure}
The full scattering process, illustrated conceptually in Fig.~\ref{fig:scattering_illu}, connects incoming and outgoing states across nine sectors defined by the number of photons and dark waves, and conserves the total lattice momentum \(\bm{P}\) and energy \(E\). We label the subspace with \(\alpha\) photons as channel \(\alpha\), for \(\alpha \in \{0, 1, 2\}\).



In our companion paper \cite{Wang2025MultiExcitation}, we demonstrate that the full, multi-channel scattering process can be constructed from the dynamics of the atomic spin waves. Because of the separation of  time  scales  between the atomic resonance and decay rate, the atomic dynamics can be described by an effective Hamiltonian \cite{Lehmberg1970,ExponentialAsenjoGarcia2017}, where photons are traced out of the system.
The scattering of atomic spin waves is thus described by a complex dispersion relation \(\epsilon^{(2)}(\bm{P}, \bm{q}) \equiv \epsilon(\bm{P}/2 + \bm{q}) + \epsilon(\bm{P}/2 - \bm{q})\). The single-excitation dispersion \(\epsilon(\bm{p})\) is purely real for dark spin waves (\(\|\bm{p}\| > p_0\)), where \(\epsilon(\bm{p}) = \Delta(\bm{p})\), but acquires an imaginary part for bright spin waves (\(\|\bm{p}\| < p_0\)), \(\epsilon(\bm{p}) = \Delta(\bm{p}) - i\Gamma(\bm{p})/2\), where $\Gamma(\bm{p})$  represents their decay rate into photons \cite{bettles2015cooperative,Adams2019}.




 \emph{Universal behavior in the S-matrix.---} We first focus on the purely atomic scattering within the subspace of two dark spin waves (channel \(\alpha=0\)). The corresponding S-matrix, \(S_0(\bm{P}, E)\), which acts on the manifold of two-excitation dark spin waves \(\mathcal{M}_0(\bm{P},E)\) with the same total momentum \(\bm{P}\) and energy \(E\), has a simple structure due to the momentum-independence of the interaction. It possesses a single non-trivial eigenvalue \(s(\bm{P}, E)\) with a corresponding normalized eigenvector of the form~\cite{Wang2025MultiExcitation}
\begin{equation}
|\psi(\bm{P}, E)\rangle = \mathcal{N} \int_{\mathcal{M}_0(\bm{P}, E)} d\bm{q} \, \frac{|\bm{P}, \bm{q}\rangle}{\sqrt{v_g(\bm{P}, \bm{q})}},
\label{eq:psi_definition}
\end{equation}
 where \(\mathcal{N}\) is a normalization factor and \(v_g(\bm{P},\bm{q})=|\nabla_{\bm{q}}\Delta^{(2)}(\bm{P},\bm{q})|\) is the relative group velocity.
The corresponding eigenvalue \(s(\bm{P}, E)\) is given by \(s(\bm{P}, E) = L(\bm{P}, E - i0) / L(\bm{P}, E + i0)\) \cite{Wang2025MultiExcitation}.
 This central quantity, the local propagator \(L(\bm{P}, \omega)\), represents the system's on-site Green's function
\eq{
L(\bm{P}, \omega) &= \int_{BZ^{(2)}} d\bm{q} \frac{1}{\omega - \epsilon^{(2)}(\bm{P}, \bm{q})}.\label{eqLint2D_final}
}
Here, the integral is over the two-excitation Brillouin zone \(BZ^{(2)}\) [Fig.~\ref{fig:results_all}(a)], which accounts for the indistinguishability of the two bosonic excitations.
Its imaginary part, \(\text{Im}[L]\), is proportional to the total density of available on-shell atomic spin-wave states (both bright and dark). Its real part represents the density of the off-shell atomic excitations.




Universality emerges as the energy \(E\) approaches the critical energies ($E_{\text{crit}}$) of the dark-spin-wave dispersion \(\Delta^{(2)}(\bm{P}, \bm{q})\). These critical energies correspond to the relative momentum \(\bm{q}\) where the group velocity vanishes, which generically occur at band maxima ($E_{\text{max}}$) and saddle points ($E_{\text{sadd}}$) [red and yellow markers in Fig.~\ref{fig:results_all}(a), respectively] \footnote{The spectrum of the dark atomic spin waves approaches $-\infty$ near the light cone, hence for square lattices there is no band minimum.} \footnote{The location and multiplicity of the critical points can be tuned by the choice of total momentum \(\bm{P}\) and the lattice geometry.}. As the energy approaches these values, the local propagator \(L\) diverges logarithmically, but the nature of this divergence is different for both types of critical energies:
\eq{
 L(\bm{P}, E\pm i0)
\propto
\begin{cases}
- \log\left(|E - E_{\text{max}}|\right) & \text{as } E \rightarrow E_{\text{max}} \\
\pm i \log\left(|E - E_{\text{sadd}}|\right) & \text{as } E \rightarrow E_{\text{sadd}}
\end{cases} \label{eqLsaddlelog_final}
}
The distinct scaling of $L$ near the two types of critical points is a direct consequence of their differing dispersion topologies.

The local dispersion is homeomorphic to \(-(q_x^2+q_y^2)\) near a maximum but to \(q_x^2-q_y^2\) near a saddle point, a difference formally captured by the transformation \(q_x \to iq_x\),  which introduces the relative factor of $i$ in the resulting integral.


This reversal of whether the real or imaginary part of \(L\) diverges leads to distinct universal limits for the S-matrix eigenvalue:
\begin{equation}
s(\bm{P}, E) \rightarrow \begin{cases}
1 & \text{as } E \rightarrow E_{\text{max}} \\[4pt]
-1 & \text{as } E \rightarrow E_{\text{sadd}}
\end{cases}
\label{eqSEmax0_final}
\end{equation}
Evaluating \(L(E\pm i0)\) numerically, we plot \(s(\bm{P}=\bm{0}, E)\) in Fig.~\ref{fig:results_all}(b), which precisely follows this universal prediction near the critical points.

At \(E \to E_{\text{max}}^-\), \(s \to 1\), meaning the S-matrix becomes the identity operator and the dark waves pass through each other without interaction. At \(E \to E_{\text{sadd}}\), \(s \to -1\), and the scattering eigenstate acquires a \(\pi\) phase shift.




The physical pictures behind these limits are distinct and intuitive. At a band maximum \(E_{\text{max}}\), the density of on-shell states is finite, but the  density of off-shell excitations
diverges (\(\text{Re}[L] \to -\infty\)). The excitations become ``dressed" by an infinitely heavy cloud of off-shell excitations, which suppresses their interaction and leads to \(s \to 1\). Conversely, at a saddle point \(E_{\text{sadd}}\), the density of available on-shell states diverges (\(|\text{Im}[L]| \to \infty\)). The incoming spin wave pair encounters this infinite reservoir of states to scatter into, which forces a strong scattering event corresponding to a \(\pi\) phase shift, and thus \(s \to -1\).





Crucially, the scattering behavior is not uniform across the equi-energy contour. At \(E = E_{\text{sadd}}\), the Hilbert space effectively fractures: the scattering eigenstate \(|\psi(\bm{P},E_{\text{sadd}})\rangle\) becomes localized only on the discrete saddle points in momentum space. This means that dark-wave pairs on the ``saddle lines" connecting these points [yellow dashed lines in Fig.~\ref{fig:results_all}(a)] are orthogonal to the scattering eigenstate and form a separate, non-interacting channel where the S-matrix is simply the identity operator.


The unitarity of the S-matrix, \(|s(\bm{P}, E)| \to 1\), at these critical energies is the central mechanism of this work, with profound consequences for the entire scattering problem.
The probability of scattering to photon-containing output channels is $1 - |s|^2$, and is non-zero for generic energies. At the critical energies, the unitarity implies that the channel of two dark spin waves becomes completely decoupled from the photon-containing channels.
This decoupling has a far-reaching effect: in all other scattering channels (i.e., for $\alpha, \beta > 0$), inter-excitation interactions are completely suppressed, and each excitation evolves independently according to single-excitation dynamics~\cite{Wang2025MultiExcitation}.




\emph{Universal features in scattering cross sections.---}
The universal S-matrix behavior has interesting and non-trivial  implications for the corresponding scattering cross sections.
When two plane-wave excitations in channel $\alpha$ impinge on each other with flux \(R\) in their center of mass frame, \(\sigma_{\alpha,\beta} R\) gives the number of pairs of excitations scattered into channel \(\beta\) per unit time.
Given an incoming state with a specific momentum, this partial scattering cross section is given by
\eq{
\sigma_{\alpha, \beta}=  \frac{4\pi a_\alpha}{v_g} |\bar{T}(\bm{P}, E+i0)|^2 \text{Im}[L_\beta(\bm{P}, E+i0)], \label{eqsigma_CS_revised}
}
where \(\bm{P}\) is the projection of the total momentum of the incoming state onto the array, and $E$ is the energy of the incoming state.
This expression is a product of four distinct factors:  (1) a prefactor $a_\alpha$, representing the conversion amplitude from the incoming state to the spin-wave modes of the array~\cite{Wang2025MultiExcitation}, which remains finite and does not contribute to the universal scaling;
(2) the inverse group velocity of the incoming state, \(v_g^{-1}\); (3) the density of states  \(\text{Im}[L_\beta]\) of $\beta$ bright spin waves and $2-\beta$ dark spin waves at $(\bm{P},E)$; and (4) the core scattering probability \(|\bar{T}(\bm{P}, E+i0)|^2\), where we introduce the atomic T-matrix   \(\bar{T}(\bm{P}, E+i0) = -L(\bm{P}, E+i0)^{-1}\).
The fact that \(|\bar{T}|^2\) is independent of the incoming and outgoing momentum is a direct consequence of the hard-core interaction, representing a uniform scattering probability into any available pair of atomic spin-wave states.
To define \(L_\beta\), we partition $\textrm{BZ}^{(2)}$ into domains \(D_\beta\) based on the number of constituent bright spin waves [see Fig.~\ref{fig:results_all}(a)].
Within this formalism, the components \(\text{Im}[L_{\beta > 0}]\) represent the density of states for the transient bright-wave modes that serve as the direct gateway to the photonic outgoing states.
This analytical transparency allows the universal properties of \(L\) to be directly mapped onto all scattering observables, making the origin of the universal phenomena particularly clear.


For
incoming states containing photons (\(\alpha=1,2\)), the group velocity \(v_g\) is finite.
 As the energy \(E\) approaches a critical energy \(E_{\text{crit}}\), \(\sigma_{\alpha,\beta}\) always vanishes logarithmically, with the specific scaling laws detailed in our summary table, Fig.~\ref{fig:summary_table}.  This universal suppression is confirmed in Fig.~\ref{fig:results_all}(d) for the case of two normally incident photons.


In stark contrast, for an incoming state of two dark spin waves (\(\alpha=0\)) with relative momentum \(\bm{q}\) approaching a critical point \(\bm{q}_{\text{crit}}\), the group velocity vanishes as \(v_g \propto \sqrt{\Delta E}\), where $\Delta E=|E-E_{\mathrm{crit}}|$. The diverging  inverse group velocity  \(v_g^{-1} \propto \Delta E^{-1/2}\) overpowers the vanishing $|T|^2\text{Im}[L]$, leading to a divergent total cross section \(\sigma_{0,\mathrm{tot}}=\sum_{\beta} \sigma_{0,\beta}\) \brown{\footnote{Divergent scattering cross sections have also been reported in spin-orbit coupled systems \cite{PhysRevA.91.022706, PhysRevA.94.022706}}}.
As also detailed in  Fig.~\ref{fig:summary_table}, the specific scaling of this divergence depends on the critical point. The behavior of the total cross section \(\sigma_{0,\text{tot}}\) as a function of the incoming relative momentum \(\bm{q}\) is plotted in  Fig.~\ref{fig:results_all}(c1).
This plot reveals that the cross section is enhanced at the critical points, which are indicated in the figure by the yellow and red markers.
Conversely, for incoming states on the saddle lines but away from the saddle points, \(v_g\) is finite, and the total cross section is consequently suppressed, as the figure also demonstrates.



 The branching ratio into channel $\beta$, $\sigma_{0,\beta}/\sigma_{0,\text{tot}}$, simplifies to $\text{Im}[L_\beta]/\text{Im}[L]$. This ratio reveals further universal signatures:
\begin{itemize}
    \item Approaching a saddle point (\(\bm{q} \to \bm{q}_{\text{sadd}}\)): The ratio of scattering into photons vanishes as the divergence of \(\text{Im}[L]\) is dominated by the dark-spin-wave contribution \(\text{Im}[L_0]\).
    \item Approaching a band maximum (\(\bm{q} \to \bm{q}_{\text{max}}\)): The ratio approaches a finite, non-zero value because all partial densities of states \(\text{Im}[L_\beta]\) remain finite.
\end{itemize}
Along the saddle lines, where the total cross section is already suppressed, the branching ratio into photons also vanishes. Figure~\ref{fig:results_all}(c2) perfectly confirms all of these distinct behaviors.







\begin{figure}
    \centering
\includegraphics[width=0.9\linewidth]{table_cs_illu3.pdf}
\caption{Universal scaling of partial cross sections \(\sigma_{\alpha,\beta}\) near critical energies (\(\Delta E = |E - E_{\text{crit}}|\)). The colors denote four scaling behaviors: \(\Delta E^{-1}\log(\Delta E)^{-1}\) (dark orange), \(\Delta E^{-1}\log(\Delta E)^{-2}\) (light orange), \(\log(\Delta E)^{-1}\) (light blue), and \(\log(\Delta E)^{-2}\) (dark blue). The tables classify the behavior for each scattering channel \((\alpha, \beta)\) at \(E_{\text{max}}\) and \(E_{\text{sadd}}\), where for \(E_{\text{sadd}}\) the \(\alpha=0\) channel is further divided based on whether the incoming relative momentum is at the saddle point (\(\bm{q}_{\text{sadd}}\)).}
    \label{fig:summary_table}
\end{figure}



Finally, we resolve an apparent paradox at $E_{\text{max}}$: the S-matrix becomes trivial ($s \to 1$), while the cross section $\sigma_{0,\text{tot}}$ diverges. Near the band edge, the trivial S-matrix implies that the scattering probability per unit time per excitation pair vanishes, but the vanishing group velocity ($v_g \to 0$) of the incoming excitations means the density of excitations within the scattering region increases (when keeping the incident flux constant). The long interaction time afforded by the slow-moving excitations enhances the total cross section, which is the rate of excitations scattered nonlinearly per unit incoming flux. In other words, the diverging cross sections reflects a subtle interplay in which the scattering probability vanishes more slowly than the interaction time grows.


















\emph{Discussion.---}
By fixing the scattering behavior at the critical energies, the universal physics discovered provides a qualitative description of the scattering behavior across a wide energy range, as illustrated in Fig.~\ref{fig:results_all} (b).
The prominent suppression of the scattering cross section is anchored precisely at the theoretically predicted \(E_{\text{sadd}}\) and \(E_{\text{max}}\), while the overall width of the interaction region is set by the natural dissipative scale of the system, \(\Gamma(\bm{p}=\bm{0})\) (here, \(\approx 6\Gamma_0\), where \(\Gamma_0\) is the single-atom decay rate). Beyond these sharp features, our results also reveal another generic property: the predominance of scattering into the dark-wave channel (\(\beta=0\)). This originates from the high density of states associated with the relatively flat dark-wave bands, which enhances the \(\sigma_{\alpha,0}\) cross section over other channels. This preferential scattering into subradiant modes is therefore a general feature of atomic arrays supporting flat collective bands~\cite{jaworowski2025laughlin}.



Experimentally, two-photon cross sections can be extracted from the linear dependence of photon transmission or reflection on input intensity \cite{Wang2025MultiExcitation}. The main experimental test is to measure this cross section as a function of energy near the critical points. Although broadening due to laser resolution, finite size~\footnote{Finite-size effects impose a minimum energy resolution, scaling as \(\Delta E_{\min} \propto 1/N^2\) for an \(N \times N\) array.},
and disorder will smooth out sharp theoretical features, observing robust suppression at the critical energies would clearly validate our theory with current capabilities.

While our analysis assumes a two-level nonlinearity and a square array, the core mechanism we identify is far more general.
 The universal phenomena arise from the singular phase space near a band-structure critical point, making the details of the short-range interaction irrelevant.
 We thus expect these universal  phenomena for a wide range of interactions and array geometries.
 The decoupling of dark spin waves from photon channels at critical energies indicates two-excitation subradiant states, discussed further in \cite{Wang2025MultiExcitation}.
This detailed two-excitation scattering analysis underpins future studies of many-body dynamics in these driven-dissipative systems.



\emph{Acknowledgement.~}
We thank Alec Douglas, Simon Hollerith and Sayed Ali Akbar Ghorashi for discussions. Y.W. and S.Y. acknowledge support from the NSF through CUA PFC (PHY-2317134), PHY-2207972, and the Q-SEnSE QLCI (OMA-2016244). O.R.-B. acknowledges support from Fundación Mauricio y Carlota Botton and from Fundació Bancaria “la Caixa” (LCF/BQ/AA18/11680093).  V. W. acknowledges support from the NSF through Grant No. PHYS-2409630 for early-stage investigations of the scattering formalism, and support from the U.S. Department of Energy, Office of Science, Office of Basic Energy Sciences Energy Frontier Research Centers program under Award Number DE-SC0025620 for the subsequent analysis of the dynamics near critical points.



\bibliographystyle{apsrev4-2}
\bibliography{library_corrected}


\end{document}
===== END FILE: Short_arxiv_submission.tex =====
